%=====================================================================================================
\section{Conclusion}
\label{sec:conclusion}
Pseudo-label fine-tuning of a compact detector on one camera improved that camera and made the detector worse than
its pretrained weights on every other camera we tested and on COCO people. A size-blind label threshold and
forgetting explain the loss. Averaging pretrained and fine-tuned weights, with the coefficient chosen on other
cameras, recovered the test-only score (\MsSelHo{} against \MsFtHo{} for the model chosen on the adaptation camera)
at no cost. The CNN student's apparent accuracy advantage came from its input size. For the VLM operator, text in the
camera view is an attack channel, and a policy layer outside the model raised the rejection of unsafe requests from
\SfGrTestRej{} to \SfPoTestRej{} without hurting legitimate ones. Next steps are the human ground truth, on-device
measurements on a Jetson Orin, adaptive red-teaming of the policy, and more cameras, including night scenes. Every
number in the paper is generated from the released result files by \texttt{paper/make\_numbers.py}.
